Un \emph{anti-triangle de Pascal} est un tableau en forme de triangle équilatéral dans lequel sont disposés des
nombres tels que, excepté pour les nombres placés sur la ligne du bas, chaque nombre soit égal à la valeur
absolue de la différence entre les deux nombres situés juste en-dessous. Par exemple, le tableau ci-dessous est
un anti-triangle de Pascal de quatre lignes qui contient chaque entier entre $1$ et $10$.
\[
\begin{array}{ccccccc}
 & & & 4 & & & \\
 & & 2 & & 6 & & \\
 & 5 & & 7 & & 1 & \\
8 & & 3 & & 10 & & 9
\end{array}
\]
Existe-t-il un anti-triangle de Pascal de $2018$ lignes qui contient tous les entiers de $1$ à
$1 + 2 + \cdots + 2018$ ?
